% AI-written, AI cross-reviewed; no human sign-off; not a claim that RH is proved
% RH Notes blueprint source excerpts for Z2.K5
% Verbatim excerpts in the notation of the parent blueprint. The labelled editions and source ranges appear below.
% See the unit YAML for the current proof status. Earlier-round passages retain their original wording.
\section*{Z2.K5. Equivariant character count}

% Parent source SHA-256: 1af64922b98ed2310744411b1f18b25cc426c97220646321471836e1a3a409d6
% blueprint/ch_P.tex lines 517 through 648; prop:P-k-count
\begin{proposition}[Divisors, units and finite subsets]\label{prop:P-k-count}
\lean{UpperTheory.V.ZK.sCount}
\uses{prop:P-k-few}
With $v=v_\ell(E)$, $g=2n'+\lfloor\log_2 h_M\rfloor+2$, where $h_M$ is
the class number, and $X=(b+v)g+c+v$, one has
\[
 |\mathcal H_q|\le B_*:=T_{2E+1}T_{2E}\ell^X.
\]
Moreover $\Sigma(M)=s$ for an integer $s\ge0$ and $\ell^s\le B_*$.
The first assertion combines \nolinkurl{sCount} with separation
and cancellation (\nolinkurl{sNum}, \path{NumThree.lean:221--241}).
The second is \nolinkurl{PositivityTheory.QuaternionShaBounds.sSha}
(\path{OrderOfOneAddPrimePowZmodPrimePow.lean:490};
\nolinkurl{SSha}, \path{ClassCounting/LemmaKOfSteps.lean:143--147}).
\end{proposition}
\begin{proof}
In additive notation take $Y=F^\times$, $I=\operatorname{Div}(\mathcal O_F)$,
with principal divisor map $\partial:Y\to I$. On these groups write
$\delta$ for the action of $\gamma^{-1}$ minus $a$ on $Y$.
On divisors the same endomorphism sends $u(P)$ to $u(\gamma P)-a u(P)$,
so it commutes with $\partial$ (\nolinkurl{IsDat}, \path{TwoLe.lean:56--61}).
Let $U=\partial^{-1}(\delta I)$. Form
\[
 W=U/(\delta Y+Y_{\rm tors}),\qquad
 Q_{\rm cl}=I/(\partial Y+\delta I),\qquad
 Q_{\rm u}=\ker\partial/(\delta\ker\partial+Y_{\rm tors}).
\]
These are \nolinkurl{Dat.W}, \nolinkurl{Dat.ClQ}, \nolinkurl{Dat.EQ}
(\path{SelmerOrderCounts/TwoLe.lean:34--47}).
Lemma~\ref{lem:P-r3-k5} proves the following arithmetic properties.
All three groups are finite,
$|W|=|Q_{\rm u}|\,|Q_{\rm cl}|$,
$|\mathcal H_q|\le\ell^{v_\ell(|Q_{\rm cl}|)}$ with $\mathcal H_q$ finite,
and $|W[\ell]|\le\ell^g$. The powered norm gives a homomorphism $f$ on
$W$ whose image can be represented by a finite family as in
Proposition~\ref{prop:P-k-few}. Its kernel is killed by
$m\ell^{b+v}$ with $(m,\ell)=1$. See \nolinkurl{sCount}
(\path{SelmerOrderCounts.lean:20--71}) for \nolinkurl{cardW},
\nolinkurl{eqHomLe}, \nolinkurl{kerF}, \nolinkurl{fewLe} and
\nolinkurl{selInj} in this application. In particular the condition
on $W[\ell]$ is a condition at the base $M$, although the group $W$ was
constructed using $F$.

The unit factor supplies $q^r$. The project takes the logarithmic
unit lattice modulo roots of unity and appends one trivial integral
coordinate to fill the constant direction. On the full real space of
archimedean places, $\gamma$ permutes the places in $r$ cycles of
length $h$ because $M$ is totally complex. On one cycle the absolute determinant
of this permutation minus $a$ times the identity is $a^h-1$
(the characteristic polynomial is $t^h-1$). On the extra constant line
it is $a-1$. Computing the index of the resulting lattice endomorphism
by its determinant gives
\[
 |Q_{\rm u}|(a-1)=(a^h-1)^r .
\]
Lemma~\ref{lem:P-r3-k5} identifies this logarithmic lattice with the
units modulo roots of unity. The determinant and division account for the two different
integer factors in this formula (\nolinkurl{cardEQ},
\path{SelmerOrderCounts/CardEQ.lean:98--128}). Since $\kappa=1+\ell^c$
and $h=\ell^{\nu-c}$,
\[
 v_\ell(a-1)=c+v,\qquad v_\ell(a^h-1)=\nu+v.
\]
Indeed raising a principal unit to an exponent prime to $\ell$ preserves
the valuation of its difference from $1$, and raising it to $\ell$
increases that valuation by one when $c\ge1$ for odd $\ell$ and $c\ge2$
for $\ell=2$. Both facts follow by expanding $(1+t)^j$ and comparing
the first term with subsequent terms. In particular $a^h-1=q\ell^v m$
for some $m$ prime to $\ell$, also when $\nu=c$ and $h=1$. Thus
\[
 v_\ell(|Q_{\rm u}|)+c+v=r(\nu+v)\ge r\nu .
\]
This is the equality used in \nolinkurl{sCount}
(\path{SelmerOrderCounts.lean:30--50,81--98}).

Let $K_f$ be the kernel of $f$. On its $\ell$-primary subgroup $m$ acts
invertibly, so that subgroup is killed by $\ell^{b+v}$. Each multiplication
by $\ell$ from $W[\ell^i]$ to $W[\ell^{i-1}]$ for $i\ge1$ has kernel of size at
most $|W[\ell]|$. Induction gives $|W[\ell^{b+v}]|\le\ell^{(b+v)g}$.
The index of $K_f$ in $W$ is $|\operatorname{im} f|$. Taking $\ell$-parts
yields
\[
 \ell^{v_\ell(|W|)}\le|\operatorname{im}f|\,\ell^{(b+v)g}.
\]
This is \nolinkurl{pCount} (\path{SelmerOrderCounts/PCount.lean:38--92}).
Multiplying the class-character bound by $q^r$ and inserting the unit
identity gives
\[
\begin{aligned}
 |\mathcal H_q|q^r
 &\le\ell^{v_\ell(|Q_{\rm cl}|)+v_\ell(|Q_{\rm u}|)+c+v}\\
 &\le |\operatorname{im}f|\,\ell^{(b+v)g+c+v}
 \le q^r T_{2E+1}T_{2E}\ell^X .
\end{aligned}
\]
Since $q^r>0$ it can be cancelled. In particular $X$ uses $c$, which was
fixed when $M$ was chosen, and $v_\ell(E)$. It contains no $\nu$.
This is the passage from \nolinkurl{sCount}
(\path{SelmerOrderCounts.lean:91--111}) to \nolinkurl{sNum}
(\path{NumThree.lean:221--241}).

Lemma~\ref{lem:P-r3-k6} supplies the passage back to cohomology. Restriction to the infinite
cyclotomic layer injects the degree-one localization kernel into the
equivariant continuous characters there. Each character in that image
lifts, at every sufficiently large $\nu$, to an equivariant finite-level
ideal class character killed by $\ell^\nu$. For the cyclic layer with
generator chosen above there is an injection from these latter characters
into $\mathcal H_q$. The finite-level pullback determines the original
infinite-level character. The identifications use
\nolinkurl{good_of}, \nolinkurl{ClInj}, \nolinkurl{LayEx},
\nolinkurl{ShaAsm_holds}
(\path{OrderOfOneAddPrimePowZmodPrimePow.lean:26--38,390--410,454--463}).
In the source the threshold is obtained by first extending an equivariant
character to a finite level, expressing it by pullback from a finite class
group there, and bounding the exponent of its finite image. Restriction
to deeper levels preserves the bound on its image. This explains why all
sufficiently large levels work.

Take a finite subset $T$ in the image of global restriction. A maximum
of its thresholds, $c$ and $2$ gives one level for all its members.
Choose one lift per member. If the chosen finite-level images in
$\mathcal H_q$ coincide, injectivity before restriction gives equality of
the lifts, and restriction gives equality of their members in $T$.
Thus $|T|\le B_*$. An infinite image would have a finite subset of size
$B_*+1$. The image is consequently finite, and applying the same bound
to the entire image gives its order at most $B_*$. A finite
$\mathbf Z_\ell$-module has order $\ell^s$ for its length $s$ (take
successive simple quotients). Restriction is injective, hence this $s$
is $\Sigma(M)$. This finite-subset argument is \nolinkurl{card_le} and
\nolinkurl{ShaAsm_holds} (\path{OrderOfOneAddPrimePowZmodPrimePow.lean:412--448,464--490}).
\emph{Status: complete modulo Proposition~\ref{prop:P-k-few}.}
\end{proof}

